\documentclass[letterpaper,12pt]{article}
\usepackage{amsmath, amssymb, geometry, titlesec, esint, listings, xcolor, graphicx, float, caption, booktabs, bm}
\usepackage[utf8]{inputenc}

\geometry{margin=1in}

%Make section titles smaller
\titleformat{\section}
    {\large\bfseries}
    {\thesection.}{0.5em}{}

\titleformat{\subsection}
    {\normalsize\bfseries}
    {\thesubsection.}{0.5em}{}

\titleformat{\subsubsection}
    {\normalsize\itshape}
    {\thesubsubsection.}{0.5em}{}

%Customize vertical spacing
\titlespacing*{\section}{0pt}{6pt}{3pt}
\titlespacing*{\subsection}{0pt}{6pt}{3pt}
\titlespacing*{\subsubsection}{0pt}{6pt}{3pt}

\setlength{\textfloatsep}{5pt} % space between floats and text
\setlength{\floatsep}{5pt}     % space between floats
\setlength{\intextsep}{2pt}    % space between in-text floats and text

%\linespread{0.97}

\setlength{\parindent}{0pt} % no indentation
% \setlength{\parskip}{4pt} % space between paragraphs

%Figure caption sizes
\captionsetup[figure]{justification=raggedright, font=footnotesize, labelfont=bf, skip=0pt}
\captionsetup[table]{justification=raggedright, font=small, labelfont=bf, skip=2pt}

\lstset{
    language=Python,
    basicstyle=\ttfamily\footnotesize, % Base font style and size
    keywordstyle=\color{blue!60!black},
    commentstyle=\color{green!50!black},
    stringstyle=\color{teal!60!black},
    numbers=left,                      % line numbers on the left
    numberstyle=\tiny\color{gray},
    stepnumber=1,                      % line numbers every 1 line
    breaklines=true,                   % Wrap long lines
    breakatwhitespace=false,           % Line break anywhere if needed
    columns=flexible,                  % Better spacing between characters
    frame=single,                      % Add a frame around the code
    tabsize=4,                         % Set tab size to 4 spaces
    showstringspaces=false,
    inputencoding=utf8,
    extendedchars=true
}

% Customize the spacing around equations
\makeatletter
\g@addto@macro\normalsize{%
  \setlength\abovedisplayskip{5pt}%
  \setlength\belowdisplayskip{5pt}%
  \setlength\abovedisplayshortskip{5pt}%
  \setlength\belowdisplayshortskip{5pt}%
}
\makeatother

\begin{document}

%---------------------------------------------------------------------------------------------------------------

%INFO

Christian DiPietrantonio \\
ME 5311: Computational Methods to Viscous Flows \\
\today

\vspace{6pt}

%---------------------------------------------------------------------------------------------------------------

%TITLE
{\large \textbf{Computer Assignment 03:} \normalsize Calculating the Development of a Turbulent Boundary Layer}

%---------------------------------------------------------------------------------------------------------------

%DESCRIPTION OF THE PROBLEM
\section{Description of the Problem}

The objective of this project is to use the code developed in Project 02 for the laminar boundary layer to calculate the development of a turbulent boundary layer using the algebraic model described below. In the inner region of the boundary layer, the Reichardt model is used, and the eddy viscocity is given by
\begin{equation}
    \nu_t = \kappa \nu \left[ \frac{y u_{\tau}}{\nu} - y^{+}_{a} \tanh\left( \frac{y u_{\tau}}{\nu y^{+}_{a}} \right) \right]
\end{equation}

Where $y^{+}_{a} = 9.7$.

In the outer region, the Clauser model is used, which assumes that the eddy viscocity reaches a constant value
\begin{equation}
    (\nu_t)_{\text{max}} = c_1 U_e \delta(x)
\end{equation}

where $\delta(x)$ is taken to be the location where $u = 0.994 U_{\infty}$. Whenever $\nu_t$ calculated from Reichardt's model exceeds $(\nu_t)_{\text{max}}$, the Clauser model is used and the derivative of $\nu_t$ is set to zero. Model constants $\kappa = 0.41$ and $c_1 = 0.001$ are used, and the friction velocity is found by taking the numerical derivative of $u$ at the wall. For consistant results with Project 02, the expressions for $\nu_t$ and their spatial derivatives must be nondimensionalized. Further objectives of this project are to:
\begin{enumerate}
    \item Calculate the turbulent boundary layer flow for the same dimensional properties used in the laminar case. However, this time using a value of $y_{\text{max}} = 50 \delta$ and 250 points in the $y$ direction. Calculate $\delta(x)$ and $c_f (x)$ for each $x$ location and confirm that the boundary layer grows like
    \begin{equation}
        \delta(x) = 0.375 x (Re_x)^{-1/5}
    \end{equation}
    \item Adjust the results to account for the transition from the assumed initial laminar velocity profile by obtaining a value for $x_0$, the apparent origin of the turbulent boundary layer.
    \item Show that the turbulent boundary layer follows the ``law of the wall.''
    \item Determine the sensitivity of the turbulence model to the choices of $\kappa$ and $c_1$.
    \item Determine what changes are necessary in the grid resolution and stretching factor to obtain a converged solution when the Reynolds number is increased by a factor of 10.
    \item Discuss the problems encountered with the implementation of this model.
\end{enumerate}

%---------------------------------------------------------------------------------------------------------------

%DESCRIPTION OF THE NUMERICAL METHOD
\section{Description of the Numerical Method}

The numerical method used is nearly identical to Project 02, the Crank-Nicolson algorithm is used to march in the $x$-direction, a $\sinh$-based stretched grid is used in $\eta$, a fourth-order central difference scheme is used for the discretization of $\eta$ (except at the first interior points, where a second-order discretization is used), and banded linear solvers are used for both $u^*$ and $v^*$ at each $s$-step. However, for this project $\nu^*$ and $\frac{\partial \nu^*}{\partial \eta}$ are no longer constants $(\nu^* = 1, \frac{\partial \nu^*}{\partial \eta} = 0)$, instead they are computed from the turbulence model at each $x$-step. Nevertheless, the coefficients $A_j$ and $B_j$ retain the same formula as in Project 02:
\begin{equation}
    A_j = \frac{v^*_j}{J^*_j} - \frac{C}{J^{*2}_j}\frac{\partial \nu^*}{\partial \eta}\bigg|_j + \frac{C\,\nu^*_j\, J^*_{\eta,j}}{J^{*3}_j}, \qquad B_j = \frac{-C\,\nu^*_j}{J^{*2}_j}
\end{equation}

The friction velocity is nondimensionalized as $u_{\tau}^* / U_{\infty}$. From the definition $u_{\tau}^2 = \nu\,(\partial u/\partial y)|_w$ and the chain rule, the follwoing relationship is obtained:
\begin{equation}
u_{\tau}^{*2} = \frac{1}{Re_\delta \, J^{*}_0} \, \frac{\partial u^*}{\partial \eta}\bigg|_0
\end{equation}

Where $Re_\delta = U_\infty \delta / \nu_\infty$ and the wall gradient is computed via a 2nd-order forward difference stencil: 
\begin{equation}
\frac{\partial u^*}{\partial \eta}\bigg|_0 = \frac{4 u^*_1 - u^*_2}{2 \Delta \eta}
\end{equation} 

The wall coordinate is expressed as $y^+ = y^* \, u_\tau^* \, Re_\delta$, and the nondimensionalized effective viscotity as $\nu_{eff}^* = 1 + \nu_t / \nu_{\infty}$. The Reichardt and Clauser models can then be nondimensionalized by dividing by $\nu_\infty$:
\begin{equation}
\frac{\nu_t}{\nu_\infty} = \kappa \left[y^+ - y^+_a \tanh \left(\frac{y^+}{y^+_a} \right)\right], \qquad \frac{\nu_{t, \max}}{\nu_\infty} = c_1 \, Re_\delta \, y^*_{0.994}
\end{equation}

Where $y^*_{0.994}$ is the $y^*$ location where $u^* = 0.994$, found by linear interpolation of the generated data. Therefore, the nondimensional effective viscocity can be expressed in the Reichardt region as
\begin{equation}
    \nu_{eff}^* = 1 + \kappa \left[y^+ - y_{a}^+ \tanh\left(\frac{y^+}{y_{a}^+}\right)\right]
\end{equation}

and in the Clauser region as:
\begin{equation}
    \nu_{eff}^* = 1 + c_1 \, Re_\delta \, y^*_{0.994}
\end{equation}

in the Clauser region. Finally, the derivtive required for $A_j$ can be derived using the chain rule for both regions:
\begin{equation}
    \frac{\partial \nu_{eff}^*}{\partial \eta} = \begin{cases}
        \kappa \, u_\tau^* \, Re_\delta \, J^* \left[1 - \mathrm{sech}^2\!\left(\dfrac{y^+}{y^+_a}\right)\right], & \text{Reichardt region} \\[6pt]
        0, & \text{Clauser region (constant } \nu_t\text{)}
    \end{cases}
\end{equation}

At the wall ($y^+ = 0$), $\mathrm{sech}^2(0) = 1$, so $d\nu^*/d\eta = 0$. Furthermore, the skin friction coefficient can be expressed in terms of the nondimensional friction velocity as $c_f = 2(u_\tau^*)^2$. The models and logic described above can now be implemented into the program written for Project 02, the result of which is displayed in Appendix~\ref{appendix-code}.
%---------------------------------------------------------------------------------------------------------------

%PRESENTATION OF RESULTS
\section{Presentation of Results}

TBD

%---------------------------------------------------------------------------------------------------------------

%DISCUSSION OF RESULTS
\section{Discussion of Results}
\subsection{General description}

TBD
- Capping y+ to avoid overflow errors

\subsection{Accuracy and stability}

TBD
%---------------------------------------------------------------------------------------------------------------

%APPENDIX A - copy of program listing
\newpage
\appendix
\section{Copy of Program Listing}
\label{appendix-code}

\lstinputlisting[
    language=Python
]{DiPietrantonio_ME5311_Project03.py}

%---------------------------------------------------------------------------------------------------------------

\end{document}
